% Options for packages loaded elsewhere
\PassOptionsToPackage{unicode}{hyperref}
\PassOptionsToPackage{hyphens}{url}
\documentclass[
]{article}
\usepackage{xcolor}
\usepackage[margin=1in]{geometry}
\usepackage{amsmath,amssymb}
\setcounter{secnumdepth}{-\maxdimen} % remove section numbering
\usepackage{iftex}
\ifPDFTeX
  \usepackage[T1]{fontenc}
  \usepackage[utf8]{inputenc}
  \usepackage{textcomp} % provide euro and other symbols
\else % if luatex or xetex
  \usepackage{unicode-math} % this also loads fontspec
  \defaultfontfeatures{Scale=MatchLowercase}
  \defaultfontfeatures[\rmfamily]{Ligatures=TeX,Scale=1}
\fi
\usepackage{lmodern}
\ifPDFTeX\else
  % xetex/luatex font selection
\fi
% Use upquote if available, for straight quotes in verbatim environments
\IfFileExists{upquote.sty}{\usepackage{upquote}}{}
\IfFileExists{microtype.sty}{% use microtype if available
  \usepackage[]{microtype}
  \UseMicrotypeSet[protrusion]{basicmath} % disable protrusion for tt fonts
}{}
\makeatletter
\@ifundefined{KOMAClassName}{% if non-KOMA class
  \IfFileExists{parskip.sty}{%
    \usepackage{parskip}
  }{% else
    \setlength{\parindent}{0pt}
    \setlength{\parskip}{6pt plus 2pt minus 1pt}}
}{% if KOMA class
  \KOMAoptions{parskip=half}}
\makeatother
\usepackage{color}
\usepackage{fancyvrb}
\newcommand{\VerbBar}{|}
\newcommand{\VERB}{\Verb[commandchars=\\\{\}]}
\DefineVerbatimEnvironment{Highlighting}{Verbatim}{commandchars=\\\{\}}
% Add ',fontsize=\small' for more characters per line
\usepackage{framed}
\definecolor{shadecolor}{RGB}{248,248,248}
\newenvironment{Shaded}{\begin{snugshade}}{\end{snugshade}}
\newcommand{\AlertTok}[1]{\textcolor[rgb]{0.94,0.16,0.16}{#1}}
\newcommand{\AnnotationTok}[1]{\textcolor[rgb]{0.56,0.35,0.01}{\textbf{\textit{#1}}}}
\newcommand{\AttributeTok}[1]{\textcolor[rgb]{0.13,0.29,0.53}{#1}}
\newcommand{\BaseNTok}[1]{\textcolor[rgb]{0.00,0.00,0.81}{#1}}
\newcommand{\BuiltInTok}[1]{#1}
\newcommand{\CharTok}[1]{\textcolor[rgb]{0.31,0.60,0.02}{#1}}
\newcommand{\CommentTok}[1]{\textcolor[rgb]{0.56,0.35,0.01}{\textit{#1}}}
\newcommand{\CommentVarTok}[1]{\textcolor[rgb]{0.56,0.35,0.01}{\textbf{\textit{#1}}}}
\newcommand{\ConstantTok}[1]{\textcolor[rgb]{0.56,0.35,0.01}{#1}}
\newcommand{\ControlFlowTok}[1]{\textcolor[rgb]{0.13,0.29,0.53}{\textbf{#1}}}
\newcommand{\DataTypeTok}[1]{\textcolor[rgb]{0.13,0.29,0.53}{#1}}
\newcommand{\DecValTok}[1]{\textcolor[rgb]{0.00,0.00,0.81}{#1}}
\newcommand{\DocumentationTok}[1]{\textcolor[rgb]{0.56,0.35,0.01}{\textbf{\textit{#1}}}}
\newcommand{\ErrorTok}[1]{\textcolor[rgb]{0.64,0.00,0.00}{\textbf{#1}}}
\newcommand{\ExtensionTok}[1]{#1}
\newcommand{\FloatTok}[1]{\textcolor[rgb]{0.00,0.00,0.81}{#1}}
\newcommand{\FunctionTok}[1]{\textcolor[rgb]{0.13,0.29,0.53}{\textbf{#1}}}
\newcommand{\ImportTok}[1]{#1}
\newcommand{\InformationTok}[1]{\textcolor[rgb]{0.56,0.35,0.01}{\textbf{\textit{#1}}}}
\newcommand{\KeywordTok}[1]{\textcolor[rgb]{0.13,0.29,0.53}{\textbf{#1}}}
\newcommand{\NormalTok}[1]{#1}
\newcommand{\OperatorTok}[1]{\textcolor[rgb]{0.81,0.36,0.00}{\textbf{#1}}}
\newcommand{\OtherTok}[1]{\textcolor[rgb]{0.56,0.35,0.01}{#1}}
\newcommand{\PreprocessorTok}[1]{\textcolor[rgb]{0.56,0.35,0.01}{\textit{#1}}}
\newcommand{\RegionMarkerTok}[1]{#1}
\newcommand{\SpecialCharTok}[1]{\textcolor[rgb]{0.81,0.36,0.00}{\textbf{#1}}}
\newcommand{\SpecialStringTok}[1]{\textcolor[rgb]{0.31,0.60,0.02}{#1}}
\newcommand{\StringTok}[1]{\textcolor[rgb]{0.31,0.60,0.02}{#1}}
\newcommand{\VariableTok}[1]{\textcolor[rgb]{0.00,0.00,0.00}{#1}}
\newcommand{\VerbatimStringTok}[1]{\textcolor[rgb]{0.31,0.60,0.02}{#1}}
\newcommand{\WarningTok}[1]{\textcolor[rgb]{0.56,0.35,0.01}{\textbf{\textit{#1}}}}
\usepackage{graphicx}
\makeatletter
\newsavebox\pandoc@box
\newcommand*\pandocbounded[1]{% scales image to fit in text height/width
  \sbox\pandoc@box{#1}%
  \Gscale@div\@tempa{\textheight}{\dimexpr\ht\pandoc@box+\dp\pandoc@box\relax}%
  \Gscale@div\@tempb{\linewidth}{\wd\pandoc@box}%
  \ifdim\@tempb\p@<\@tempa\p@\let\@tempa\@tempb\fi% select the smaller of both
  \ifdim\@tempa\p@<\p@\scalebox{\@tempa}{\usebox\pandoc@box}%
  \else\usebox{\pandoc@box}%
  \fi%
}
% Set default figure placement to htbp
\def\fps@figure{htbp}
\makeatother
\setlength{\emergencystretch}{3em} % prevent overfull lines
\providecommand{\tightlist}{%
  \setlength{\itemsep}{0pt}\setlength{\parskip}{0pt}}
\usepackage{bookmark}
\IfFileExists{xurl.sty}{\usepackage{xurl}}{} % add URL line breaks if available
\urlstyle{same}
\hypersetup{
  pdftitle={R\_assignment\_Petersen},
  hidelinks,
  pdfcreator={LaTeX via pandoc}}

\title{R\_assignment\_Petersen}
\author{}
\date{\vspace{-2.5em}2026-03-16}

\begin{document}
\maketitle

\section{Data Inspection}\label{data-inspection}

\subsubsection{To start I imported both the fang\_et\_al\_genotypes and
snp\_position into
R}\label{to-start-i-imported-both-the-fang_et_al_genotypes-and-snp_position-into-r}

\begin{Shaded}
\begin{Highlighting}[]
\NormalTok{fang\_et\_al\_genotype\_URL }\OtherTok{\textless{}{-}} \StringTok{"https://raw.githubusercontent.com/EEOB{-}BioData/BCB546\_Spring2024/main/assignments/UNIX\_Assignment/fang\_et\_al\_genotypes.txt"}
\NormalTok{snp\_position\_URL }\OtherTok{\textless{}{-}}\StringTok{"https://raw.githubusercontent.com/EEOB{-}BioData/BCB546\_Spring2024/main/assignments/UNIX\_Assignment/snp\_position.txt"}
\NormalTok{fang\_et\_al\_genotypes }\OtherTok{\textless{}{-}} \FunctionTok{read.delim}\NormalTok{(fang\_et\_al\_genotype\_URL, }\AttributeTok{header =} \ConstantTok{TRUE}\NormalTok{, }\AttributeTok{sep =} \StringTok{"}\SpecialCharTok{\textbackslash{}t}\StringTok{"}\NormalTok{, }\AttributeTok{quote =} \StringTok{""}\NormalTok{)}
\NormalTok{snp\_position}\OtherTok{\textless{}{-}} \FunctionTok{read.delim}\NormalTok{(snp\_position\_URL, }\AttributeTok{header =} \ConstantTok{TRUE}\NormalTok{, }\AttributeTok{sep =} \StringTok{"}\SpecialCharTok{\textbackslash{}t}\StringTok{"}\NormalTok{, }\AttributeTok{quote =} \StringTok{""}\NormalTok{)}
\end{Highlighting}
\end{Shaded}

\begin{Shaded}
\begin{Highlighting}[]
\FunctionTok{library}\NormalTok{(tidyverse) }\CommentTok{\# making sure tidyverse is on}
\end{Highlighting}
\end{Shaded}

\begin{verbatim}
## -- Attaching core tidyverse packages ------------------------ tidyverse 2.0.0 --
## v dplyr     1.1.4     v readr     2.2.0
## v forcats   1.0.1     v stringr   1.5.2
## v ggplot2   4.0.2     v tibble    3.3.0
## v lubridate 1.9.5     v tidyr     1.3.1
## v purrr     1.1.0     
## -- Conflicts ------------------------------------------ tidyverse_conflicts() --
## x dplyr::filter() masks stats::filter()
## x dplyr::lag()    masks stats::lag()
## i Use the conflicted package (<http://conflicted.r-lib.org/>) to force all conflicts to become errors
\end{verbatim}

\subsubsection{Type of tells you the type of vector the data is. In this
case it is list for both data
files}\label{type-of-tells-you-the-type-of-vector-the-data-is.-in-this-case-it-is-list-for-both-data-files}

\begin{Shaded}
\begin{Highlighting}[]
\FunctionTok{typeof}\NormalTok{(fang\_et\_al\_genotypes) }\CommentTok{\#List}
\end{Highlighting}
\end{Shaded}

\begin{verbatim}
## [1] "list"
\end{verbatim}

\begin{Shaded}
\begin{Highlighting}[]
\FunctionTok{typeof}\NormalTok{(snp\_position) }\CommentTok{\#List}
\end{Highlighting}
\end{Shaded}

\begin{verbatim}
## [1] "list"
\end{verbatim}

\subsubsection{To determine the type of list we can use the class
function}\label{to-determine-the-type-of-list-we-can-use-the-class-function}

\begin{Shaded}
\begin{Highlighting}[]
\FunctionTok{class}\NormalTok{(fang\_et\_al\_genotypes) }\CommentTok{\# data frame}
\end{Highlighting}
\end{Shaded}

\begin{verbatim}
## [1] "data.frame"
\end{verbatim}

\begin{Shaded}
\begin{Highlighting}[]
\FunctionTok{class}\NormalTok{(snp\_position) }\CommentTok{\# data frame}
\end{Highlighting}
\end{Shaded}

\begin{verbatim}
## [1] "data.frame"
\end{verbatim}

\subsubsection{Length shows the length of the data
file}\label{length-shows-the-length-of-the-data-file}

\begin{Shaded}
\begin{Highlighting}[]
\FunctionTok{length}\NormalTok{(fang\_et\_al\_genotypes) }\CommentTok{\# 986}
\end{Highlighting}
\end{Shaded}

\begin{verbatim}
## [1] 986
\end{verbatim}

\begin{Shaded}
\begin{Highlighting}[]
\FunctionTok{length}\NormalTok{(snp\_position) }\CommentTok{\# 15}
\end{Highlighting}
\end{Shaded}

\begin{verbatim}
## [1] 15
\end{verbatim}

\subsubsection{dim shows the dimensions of the data frame: number of
rows and number of
columns}\label{dim-shows-the-dimensions-of-the-data-frame-number-of-rows-and-number-of-columns}

\begin{Shaded}
\begin{Highlighting}[]
\FunctionTok{dim}\NormalTok{(fang\_et\_al\_genotypes) }\CommentTok{\# 2782 986}
\end{Highlighting}
\end{Shaded}

\begin{verbatim}
## [1] 2782  986
\end{verbatim}

\begin{Shaded}
\begin{Highlighting}[]
\FunctionTok{dim}\NormalTok{(snp\_position) }\CommentTok{\# 983 15}
\end{Highlighting}
\end{Shaded}

\begin{verbatim}
## [1] 983  15
\end{verbatim}

\subsubsection{The names of the columns are shown below only for the
snp\_positions.}\label{the-names-of-the-columns-are-shown-below-only-for-the-snp_positions.}

\#\#\#Because there were 986 column names I did not list them all here
but the first three are listed below \#\#\# - Sample\_ID, \#\#\# -
JG\_OTU, \#\#\# - Group

\begin{Shaded}
\begin{Highlighting}[]
\FunctionTok{colnames}\NormalTok{(snp\_position) }\CommentTok{\# [1] "SNP\_ID"               "cdv\_marker\_id"        "Chromosome"           "Position"             "alt\_pos"             }
                       \CommentTok{\# [6] "mult\_positions"       "amplicon"             "cdv\_map\_feature.name" "gene"                 "candidate.random"    }
                       \CommentTok{\# [11] "Genaissance\_daa\_id"   "Sequenom\_daa\_id"      "count\_amplicons"      "count\_cmf"            "count\_gene"          }
\FunctionTok{colnames}\NormalTok{(fang\_et\_al\_genotypes)}
\end{Highlighting}
\end{Shaded}

\subsubsection{The size of the files}\label{the-size-of-the-files}

\begin{Shaded}
\begin{Highlighting}[]
\FunctionTok{object.size}\NormalTok{(fang\_et\_al\_genotypes) }\CommentTok{\# 22681376 bytes}
\end{Highlighting}
\end{Shaded}

\begin{verbatim}
## 22681376 bytes
\end{verbatim}

\begin{Shaded}
\begin{Highlighting}[]
\FunctionTok{object.size}\NormalTok{(snp\_position) }\CommentTok{\# 327392 bytes}
\end{Highlighting}
\end{Shaded}

\begin{verbatim}
## 327392 bytes
\end{verbatim}

\subsubsection{I used the command glimpse instead of str to look at the
data because the glimpse output is cleaner and easier to
read.}\label{i-used-the-command-glimpse-instead-of-str-to-look-at-the-data-because-the-glimpse-output-is-cleaner-and-easier-to-read.}

\subsubsection{This shows a neat view of the data in the
files}\label{this-shows-a-neat-view-of-the-data-in-the-files}

\subsubsection{\texorpdfstring{\textbf{To observe the data delete
eval=FALSE}}{To observe the data delete eval=FALSE}}\label{to-observe-the-data-delete-evalfalse}

\begin{Shaded}
\begin{Highlighting}[]
\FunctionTok{glimpse}\NormalTok{(fang\_et\_al\_genotypes)}
\FunctionTok{glimpse}\NormalTok{(snp\_position)}
\end{Highlighting}
\end{Shaded}

\subsubsection{Also because we are working in R, once the data has been
uploaded it is easy to look at the full data frame in table using the
view
command}\label{also-because-we-are-working-in-r-once-the-data-has-been-uploaded-it-is-easy-to-look-at-the-full-data-frame-in-table-using-the-view-command}

\begin{Shaded}
\begin{Highlighting}[]
\FunctionTok{view}\NormalTok{(fang\_et\_al\_genotypes)}
\FunctionTok{view}\NormalTok{(snp\_position)}
\end{Highlighting}
\end{Shaded}

\subsection{This also helps with understanding what the variables in the
data sets. Below I have described the variables for each of the data
sets that is important for the
assignment}\label{this-also-helps-with-understanding-what-the-variables-in-the-data-sets.-below-i-have-described-the-variables-for-each-of-the-data-sets-that-is-important-for-the-assignment}

\subsubsection{fang\_et\_al\_genotypes:}\label{fang_et_al_genotypes}

\subsubsection{- Group --\textgreater{} This tells you what organism
that data came from. for our assignment we are intersted
in:}\label{group-this-tells-you-what-organism-that-data-came-from.-for-our-assignment-we-are-intersted-in}

\subsubsection{- Maize group --\textgreater{} ZMMIL, ZMMLR, and
ZMMMR}\label{maize-group-zmmil-zmmlr-and-zmmmr}

\subsubsection{- Teosinte group --\textgreater{} ZMPBA, ZMPIL, and
ZMPJA}\label{teosinte-group-zmpba-zmpil-and-zmpja}

\subsubsection{- The rest of the column titles are the SNP\_IDs with the
genotype for that specific
snp}\label{the-rest-of-the-column-titles-are-the-snp_ids-with-the-genotype-for-that-specific-snp}

\subsubsection{snp\_position:}\label{snp_position}

\subsubsection{- SNP\_ID --\textgreater{} This is the snp
identifier}\label{snp_id-this-is-the-snp-identifier}

\subsubsection{- Chromosome --\textgreater{} This is the chromosome on
which the SNP is
located}\label{chromosome-this-is-the-chromosome-on-which-the-snp-is-located}

\subsubsection{- Position --\textgreater{} This is the location (bp) of
the SNP on that
chromosome}\label{position-this-is-the-location-bp-of-the-snp-on-that-chromosome}

\section{Data Processing}\label{data-processing}

\subsubsection{First I wanted to pull out the Maize and Teosinte groups
from the fang\_et\_al\_genotypes data
frame}\label{first-i-wanted-to-pull-out-the-maize-and-teosinte-groups-from-the-fang_et_al_genotypes-data-frame}

\subsubsection{I created new vectors titled Maize genotypes that
contained the groups ZMMIL, ZMMLR, and ZMMMR and Teosinte\_genotypes
that contained the groups ZMPBA, ZMPIL, and
ZMPJA}\label{i-created-new-vectors-titled-maize-genotypes-that-contained-the-groups-zmmil-zmmlr-and-zmmmr-and-teosinte_genotypes-that-contained-the-groups-zmpba-zmpil-and-zmpja}

\begin{Shaded}
\begin{Highlighting}[]
\NormalTok{Maize\_genotype }\OtherTok{\textless{}{-}}\NormalTok{ fang\_et\_al\_genotypes }\SpecialCharTok{\%\textgreater{}\%} \FunctionTok{filter}\NormalTok{(Group }\SpecialCharTok{\%in\%} \FunctionTok{c}\NormalTok{(}\StringTok{"ZMMIL"}\NormalTok{, }\StringTok{"ZMMLR"}\NormalTok{, }\StringTok{"ZMMMR"}\NormalTok{))}
\NormalTok{Teosinte\_genotype }\OtherTok{\textless{}{-}}\NormalTok{ fang\_et\_al\_genotypes }\SpecialCharTok{\%\textgreater{}\%} \FunctionTok{filter}\NormalTok{(Group }\SpecialCharTok{\%in\%} \FunctionTok{c}\NormalTok{(}\StringTok{"ZMPBA"}\NormalTok{, }\StringTok{"ZMPIL"}\NormalTok{, }\StringTok{"ZMPJA"}\NormalTok{))}
\end{Highlighting}
\end{Shaded}

\subsubsection{Next I wanted to transpose the mazie and teosinte vectors
so they could be joined with the SNP\_positions
data}\label{next-i-wanted-to-transpose-the-mazie-and-teosinte-vectors-so-they-could-be-joined-with-the-snp_positions-data}

\subsubsection{This created two new files, Mazie\_genotype\_tran and
Teosinte\_genotype\_tran}\label{this-created-two-new-files-mazie_genotype_tran-and-teosinte_genotype_tran}

\begin{Shaded}
\begin{Highlighting}[]
\NormalTok{Maize\_genotype\_tran }\OtherTok{\textless{}{-}} \FunctionTok{t}\NormalTok{(Maize\_genotype)}
\NormalTok{Teosinte\_genotype\_tran }\OtherTok{\textless{}{-}} \FunctionTok{t}\NormalTok{(Teosinte\_genotype)}
\end{Highlighting}
\end{Shaded}

\subsubsection{I was running into issues with manipulating the
Maize\_genotype\_tran and
Teosinte\_genotype\_tran}\label{i-was-running-into-issues-with-manipulating-the-maize_genotype_tran-and-teosinte_genotype_tran}

\subsubsection{I used the the typeof() function to determine what type
of vector my new data
was}\label{i-used-the-the-typeof-function-to-determine-what-type-of-vector-my-new-data-was}

\subsubsection{Using these commands I determined that after using the
transpose function the data is converted to a character
vector}\label{using-these-commands-i-determined-that-after-using-the-transpose-function-the-data-is-converted-to-a-character-vector}

\begin{Shaded}
\begin{Highlighting}[]
\FunctionTok{typeof}\NormalTok{(Maize\_genotype\_tran)}
\end{Highlighting}
\end{Shaded}

\begin{verbatim}
## [1] "character"
\end{verbatim}

\begin{Shaded}
\begin{Highlighting}[]
\FunctionTok{typeof}\NormalTok{(Teosinte\_genotype\_tran)}
\end{Highlighting}
\end{Shaded}

\begin{verbatim}
## [1] "character"
\end{verbatim}

\subsubsection{I used the as.data.frame() function to change my vectors
back to data frames for easier
manipulation}\label{i-used-the-as.data.frame-function-to-change-my-vectors-back-to-data-frames-for-easier-manipulation}

\begin{Shaded}
\begin{Highlighting}[]
\NormalTok{Maize\_genotype\_tran }\OtherTok{\textless{}{-}} \FunctionTok{as.data.frame}\NormalTok{(Maize\_genotype\_tran)}
\NormalTok{Teosinte\_genotype\_tran }\OtherTok{\textless{}{-}} \FunctionTok{as.data.frame}\NormalTok{(Teosinte\_genotype\_tran)}
\end{Highlighting}
\end{Shaded}

\subsubsection{These steps are to fix the column names so they can be
merged with the snp\_positions
file}\label{these-steps-are-to-fix-the-column-names-so-they-can-be-merged-with-the-snp_positions-file}

\begin{Shaded}
\begin{Highlighting}[]
\FunctionTok{colnames}\NormalTok{(Maize\_genotype\_tran) }\OtherTok{\textless{}{-}}\NormalTok{ Maize\_genotype\_tran[}\DecValTok{1}\NormalTok{,] }\CommentTok{\# Takes the data in the first row and makes that the column names}
\FunctionTok{colnames}\NormalTok{(Teosinte\_genotype\_tran) }\OtherTok{\textless{}{-}}\NormalTok{ Teosinte\_genotype\_tran[}\DecValTok{1}\NormalTok{,] }\CommentTok{\# Takes the data in the first row and makes that the column names}
\NormalTok{Maize\_genotype\_tran\_2 }\OtherTok{\textless{}{-}}\NormalTok{ Maize\_genotype\_tran[}\SpecialCharTok{{-}}\NormalTok{(}\DecValTok{2}\SpecialCharTok{:}\DecValTok{3}\NormalTok{),] }\CommentTok{\# Deleting rows 2\&3}
\NormalTok{Teosinte\_genotype\_tran\_2 }\OtherTok{\textless{}{-}}\NormalTok{ Teosinte\_genotype\_tran[}\SpecialCharTok{{-}}\NormalTok{(}\DecValTok{2}\SpecialCharTok{:}\DecValTok{3}\NormalTok{),] }\CommentTok{\# Deleting rows 2\&3}
\end{Highlighting}
\end{Shaded}

\subsubsection{Now I need to adjust the row names for the maize and
teosinte data frames so Sample\_ID is a row name with the correct data
in the column. This will be needed to merge the Maize and Teosinte data
with the snp\_position
data}\label{now-i-need-to-adjust-the-row-names-for-the-maize-and-teosinte-data-frames-so-sample_id-is-a-row-name-with-the-correct-data-in-the-column.-this-will-be-needed-to-merge-the-maize-and-teosinte-data-with-the-snp_position-data}

\subsubsection{I used the mutate function to add a new column to the
data frame. This moved the sample\_ID names from being row titles to
being data in the data
frame}\label{i-used-the-mutate-function-to-add-a-new-column-to-the-data-frame.-this-moved-the-sample_id-names-from-being-row-titles-to-being-data-in-the-data-frame}

\begin{Shaded}
\begin{Highlighting}[]
\NormalTok{Maize\_genotype\_tran\_3 }\OtherTok{\textless{}{-}}\NormalTok{ Maize\_genotype\_tran\_2 }\SpecialCharTok{\%\textgreater{}\%} \FunctionTok{mutate}\NormalTok{(}\AttributeTok{SAMPLE\_ID =} \FunctionTok{rownames}\NormalTok{(Maize\_genotype\_tran\_2))}
\NormalTok{Maize\_genotype\_tran\_3 }\OtherTok{\textless{}{-}} \FunctionTok{select}\NormalTok{(Maize\_genotype\_tran\_3, SAMPLE\_ID, }\FunctionTok{everything}\NormalTok{() ) }\CommentTok{\# Moving the sample\_ID to the front}
\NormalTok{Teosinte\_genotype\_tran\_3 }\OtherTok{\textless{}{-}}\NormalTok{ Teosinte\_genotype\_tran\_2 }\SpecialCharTok{\%\textgreater{}\%} \FunctionTok{mutate}\NormalTok{(}\AttributeTok{SAMPLE\_ID =} \FunctionTok{rownames}\NormalTok{(Teosinte\_genotype\_tran\_2))}
\NormalTok{Teosinte\_genotype\_tran\_3 }\OtherTok{\textless{}{-}} \FunctionTok{select}\NormalTok{(Teosinte\_genotype\_tran\_3, SAMPLE\_ID, }\FunctionTok{everything}\NormalTok{()) }\CommentTok{\# Moving the sample\_ID to the front}
\end{Highlighting}
\end{Shaded}

\subsubsection{Cutting down the snp\_position files to have only the
data I
need}\label{cutting-down-the-snp_position-files-to-have-only-the-data-i-need}

\begin{Shaded}
\begin{Highlighting}[]
\NormalTok{snp\_position\_2 }\OtherTok{\textless{}{-}}\NormalTok{ snp\_position[ ,}\SpecialCharTok{{-}}\NormalTok{(}\DecValTok{5}\SpecialCharTok{:}\DecValTok{15}\NormalTok{)]}
\end{Highlighting}
\end{Shaded}

\subsubsection{I used the Merge function to merge the snp\_position data
with the Maize and Teosinte
data.}\label{i-used-the-merge-function-to-merge-the-snp_position-data-with-the-maize-and-teosinte-data.}

\subsubsection{I merged the data by the sample\_ID column from the Maize
data and the snp\_ID column from the snp\_positions
data.}\label{i-merged-the-data-by-the-sample_id-column-from-the-maize-data-and-the-snp_id-column-from-the-snp_positions-data.}

\subsubsection{I also made sure to keep the other columns in the
snp\_positions data because I need the chromosome number later
on}\label{i-also-made-sure-to-keep-the-other-columns-in-the-snp_positions-data-because-i-need-the-chromosome-number-later-on}

\begin{Shaded}
\begin{Highlighting}[]
\NormalTok{Maize\_merge }\OtherTok{\textless{}{-}} \FunctionTok{merge}\NormalTok{(Maize\_genotype\_tran\_3, snp\_position\_2, }\AttributeTok{by.x=} \StringTok{"SAMPLE\_ID"}\NormalTok{, }\AttributeTok{by.y=} \StringTok{"SNP\_ID"}\NormalTok{, }\AttributeTok{all.y=}\ConstantTok{TRUE}\NormalTok{)}
\NormalTok{Maize\_merge\_2 }\OtherTok{\textless{}{-}} \FunctionTok{select}\NormalTok{(Maize\_merge, SAMPLE\_ID, Chromosome, Position, }\FunctionTok{everything}\NormalTok{()) }\CommentTok{\# Moving the Chromosome and Position number to the front}
\NormalTok{Teosinte\_merge }\OtherTok{\textless{}{-}} \FunctionTok{merge}\NormalTok{(Teosinte\_genotype\_tran\_3, snp\_position\_2, }\AttributeTok{by.x=} \StringTok{"SAMPLE\_ID"}\NormalTok{, }\AttributeTok{by.y=} \StringTok{"SNP\_ID"}\NormalTok{, }\AttributeTok{all.y=}\ConstantTok{TRUE}\NormalTok{)}
\NormalTok{Teosinte\_merge\_2 }\OtherTok{\textless{}{-}} \FunctionTok{select}\NormalTok{(Teosinte\_merge, SAMPLE\_ID, Chromosome, Position, }\FunctionTok{everything}\NormalTok{()) }\CommentTok{\# Moving the Chromosome and Position number to the front}
\end{Highlighting}
\end{Shaded}

\subsubsection{To create a function to iterate over all the chromosomes
I wanted to see if both maize and teosinte had the same number of unique
chromosomes}\label{to-create-a-function-to-iterate-over-all-the-chromosomes-i-wanted-to-see-if-both-maize-and-teosinte-had-the-same-number-of-unique-chromosomes}

\subsubsection{If the vectors weren't the same length then the iteration
in the next step wouldn't
work}\label{if-the-vectors-werent-the-same-length-then-the-iteration-in-the-next-step-wouldnt-work}

\subsubsection{I created a vector with all the unique maize and teosinte
chromosomes}\label{i-created-a-vector-with-all-the-unique-maize-and-teosinte-chromosomes}

\begin{Shaded}
\begin{Highlighting}[]
\NormalTok{chr\_maize }\OtherTok{\textless{}{-}} \FunctionTok{unique}\NormalTok{(Maize\_merge\_2}\SpecialCharTok{$}\NormalTok{Chromosome)  }
\NormalTok{chr\_teosinte }\OtherTok{\textless{}{-}} \FunctionTok{unique}\NormalTok{(Teosinte\_merge\_2}\SpecialCharTok{$}\NormalTok{Chromosome) }
\end{Highlighting}
\end{Shaded}

\subsubsection{I am using an if else statement to compare the lengths of
the
vectors.}\label{i-am-using-an-if-else-statement-to-compare-the-lengths-of-the-vectors.}

\subsubsection{The addition of the sort function orders the chromosomes
so they can be compared evenly. !identical means not
identical}\label{the-addition-of-the-sort-function-orders-the-chromosomes-so-they-can-be-compared-evenly.-identical-means-not-identical}

\subsubsection{``Maize and Teosinte vectors are equal'' was printed so I
can continue to the function and iteration
steps}\label{maize-and-teosinte-vectors-are-equal-was-printed-so-i-can-continue-to-the-function-and-iteration-steps}

\begin{Shaded}
\begin{Highlighting}[]
\ControlFlowTok{if}\NormalTok{ (}\SpecialCharTok{!}\FunctionTok{identical}\NormalTok{(}\FunctionTok{sort}\NormalTok{(chr\_maize), }\FunctionTok{sort}\NormalTok{(chr\_teosinte))) \{}
  \FunctionTok{print}\NormalTok{(}\StringTok{"Maize and Teosinte vectors are not equal... STOP"}\NormalTok{) }
\NormalTok{\} }\ControlFlowTok{else}\NormalTok{ \{}
   \FunctionTok{print}\NormalTok{(}\StringTok{"Maize and Teosinte vectors are equal"}\NormalTok{)}
\NormalTok{  \}}
\end{Highlighting}
\end{Shaded}

\begin{verbatim}
## [1] "Maize and Teosinte vectors are equal"
\end{verbatim}

\subsubsection{Writing a function to be able to iterate the formation of
the 40 necessary
files}\label{writing-a-function-to-be-able-to-iterate-the-formation-of-the-40-necessary-files}

\subsubsection{Create\_files is the main function that I will use for
lapply further
down.}\label{create_files-is-the-main-function-that-i-will-use-for-lapply-further-down.}

\subsubsection{Iterate\_chr is the function that is doing the majority
of the work for the iteration
process}\label{iterate_chr-is-the-function-that-is-doing-the-majority-of-the-work-for-the-iteration-process}

\begin{Shaded}
\begin{Highlighting}[]
\NormalTok{Create\_files }\OtherTok{\textless{}{-}} \ControlFlowTok{function}\NormalTok{(chr) \{ }\CommentTok{\# I am using chr as the argument so the overall function will iterate over each chromosome}
\NormalTok{  Iterate\_chr }\OtherTok{\textless{}{-}} \ControlFlowTok{function}\NormalTok{(sorted\_data, }\AttributeTok{desc =} \ConstantTok{FALSE}\NormalTok{) \{ }\CommentTok{\# Here I am using sorted\_data as the argument that is defined int eh following lines}
\NormalTok{    sorted\_data }\SpecialCharTok{\%\textgreater{}\%}
      \FunctionTok{filter}\NormalTok{(Chromosome }\SpecialCharTok{==}\NormalTok{ chr) }\SpecialCharTok{\%\textgreater{}\%}
      \FunctionTok{arrange}\NormalTok{(}\ControlFlowTok{if}\NormalTok{(desc) }\FunctionTok{desc}\NormalTok{(Position) }\ControlFlowTok{else}\NormalTok{ Position) }\SpecialCharTok{\%\textgreater{}\%}
      \FunctionTok{mutate}\NormalTok{(}\ControlFlowTok{if}\NormalTok{(desc) }\FunctionTok{across}\NormalTok{(}\FunctionTok{everything}\NormalTok{(), }\SpecialCharTok{\textasciitilde{}}\FunctionTok{str\_replace\_all}\NormalTok{(}\FunctionTok{as.character}\NormalTok{(.), }\StringTok{"}\SpecialCharTok{\textbackslash{}\textbackslash{}}\StringTok{?"}\NormalTok{, }\StringTok{"{-}"}\NormalTok{)))}
\NormalTok{  \}}
\NormalTok{  maize\_inc }\OtherTok{\textless{}{-}} \FunctionTok{Iterate\_chr}\NormalTok{(Maize\_merge\_2, }\ConstantTok{FALSE}\NormalTok{)}
\NormalTok{  maize\_dec }\OtherTok{\textless{}{-}} \FunctionTok{Iterate\_chr}\NormalTok{(Maize\_merge\_2, }\ConstantTok{TRUE}\NormalTok{)}
\NormalTok{  teosinte\_inc }\OtherTok{\textless{}{-}} \FunctionTok{Iterate\_chr}\NormalTok{(Teosinte\_merge\_2, }\ConstantTok{FALSE}\NormalTok{)}
\NormalTok{  teosinte\_dec }\OtherTok{\textless{}{-}} \FunctionTok{Iterate\_chr}\NormalTok{(Teosinte\_merge\_2, }\ConstantTok{TRUE}\NormalTok{)}
  
  \FunctionTok{write\_csv}\NormalTok{(maize\_inc, }\FunctionTok{paste0}\NormalTok{(}\StringTok{"maize\_chr\_"}\NormalTok{, chr, }\StringTok{"\_increase.csv"}\NormalTok{))}
  \FunctionTok{write\_csv}\NormalTok{(maize\_dec, }\FunctionTok{paste0}\NormalTok{(}\StringTok{"maize\_chr\_"}\NormalTok{, chr, }\StringTok{"\_decrease.csv"}\NormalTok{))}
  \FunctionTok{write\_csv}\NormalTok{(teosinte\_inc, }\FunctionTok{paste0}\NormalTok{(}\StringTok{"teosinte\_chr\_"}\NormalTok{, chr, }\StringTok{"\_increase.csv"}\NormalTok{))}
  \FunctionTok{write\_csv}\NormalTok{(teosinte\_dec, }\FunctionTok{paste0}\NormalTok{(}\StringTok{"teosinte\_chr\_"}\NormalTok{, chr, }\StringTok{"\_decrease.csv"}\NormalTok{))}
\NormalTok{\}}
\end{Highlighting}
\end{Shaded}

\subsubsection{The data\_sorted is going to filter out only the
chromosome column that matches the chr that is currently being iterated
over}\label{the-data_sorted-is-going-to-filter-out-only-the-chromosome-column-that-matches-the-chr-that-is-currently-being-iterated-over}

\subsubsection{Then it is going to arrange by position by ascending or
descending values. I used the desc=FALSE to aid in the
process}\label{then-it-is-going-to-arrange-by-position-by-ascending-or-descending-values.-i-used-the-descfalse-to-aid-in-the-process}

\subsubsection{Essentially if the logic returns a TRUE then it will
sorted in descending order. If the logic returns a FALSE then it will
sort in ascending
order}\label{essentially-if-the-logic-returns-a-true-then-it-will-sorted-in-descending-order.-if-the-logic-returns-a-false-then-it-will-sort-in-ascending-order}

\subsubsection{Then using the mutate command and across(everything())
command I was able to change the ? to - for the decreasing position
files}\label{then-using-the-mutate-command-and-acrosseverything-command-i-was-able-to-change-the-to---for-the-decreasing-position-files}

\subsubsection{The \textbackslash? allows the ? to act like a character
and not a ``special character'' in
R}\label{the-allows-the-to-act-like-a-character-and-not-a-special-character-in-r}

\subsubsection{Now I am using the Iterate\_chr function to create 4
different categories of the sorted
data}\label{now-i-am-using-the-iterate_chr-function-to-create-4-different-categories-of-the-sorted-data}

\subsubsection{Using the write\_csv command I am taking the four file
versions created above and saving them as each of the 40 files needed
with the chr number changing depending on which chromosome is being
worked on by the
function}\label{using-the-write_csv-command-i-am-taking-the-four-file-versions-created-above-and-saving-them-as-each-of-the-40-files-needed-with-the-chr-number-changing-depending-on-which-chromosome-is-being-worked-on-by-the-function}

\subsubsection{The paste0 command just combines the naming convention in
the
parenthesis}\label{the-paste0-command-just-combines-the-naming-convention-in-the-parenthesis}

\subsubsection{Finally I used lapply to run Create\_files for all the
chromosomes. Using the invisible command prevents lapply from printing
the results on to the console. It keeps things looking
cleaner}\label{finally-i-used-lapply-to-run-create_files-for-all-the-chromosomes.-using-the-invisible-command-prevents-lapply-from-printing-the-results-on-to-the-console.-it-keeps-things-looking-cleaner}

\subsubsection{\texorpdfstring{\textbf{The output of the lapply is the
40 individual CSV
files}}{The output of the lapply is the 40 individual CSV files}}\label{the-output-of-the-lapply-is-the-40-individual-csv-files}

\begin{Shaded}
\begin{Highlighting}[]
\FunctionTok{invisible}\NormalTok{(}\FunctionTok{lapply}\NormalTok{(}\DecValTok{1}\SpecialCharTok{:}\DecValTok{10}\NormalTok{, Create\_files))}
\end{Highlighting}
\end{Shaded}

\section{Data Visualization}\label{data-visualization}

\subsubsection{Putting the data tables in the long format instead of the
wide format helped To make it tidy and easier to
plot}\label{putting-the-data-tables-in-the-long-format-instead-of-the-wide-format-helped-to-make-it-tidy-and-easier-to-plot}

\subsubsection{Using the pivot\_longer command it flipped the data so
all of the sample\_IDs are in one column with their respective SNP and
genotypes}\label{using-the-pivot_longer-command-it-flipped-the-data-so-all-of-the-sample_ids-are-in-one-column-with-their-respective-snp-and-genotypes}

\subsubsection{I has to use the mutate(across(everything)),
as.character() to tell R to treat the entire data table as a character
because previosuly there was a mix of characters and integers and R
doesn't like
that}\label{i-has-to-use-the-mutateacrosseverything-as.character-to-tell-r-to-treat-the-entire-data-table-as-a-character-because-previosuly-there-was-a-mix-of-characters-and-integers-and-r-doesnt-like-that}

\subsubsection{I also kept the the chromosome and position columns when
flipping the data
table}\label{i-also-kept-the-the-chromosome-and-position-columns-when-flipping-the-data-table}

\begin{Shaded}
\begin{Highlighting}[]
\NormalTok{Maize\_long }\OtherTok{\textless{}{-}}\NormalTok{ Maize\_merge\_2 }\SpecialCharTok{\%\textgreater{}\%} \FunctionTok{mutate}\NormalTok{(}\FunctionTok{across}\NormalTok{(}\FunctionTok{everything}\NormalTok{(), as.character)) }\SpecialCharTok{\%\textgreater{}\%} 
  \FunctionTok{pivot\_longer}\NormalTok{(}\AttributeTok{cols =} \SpecialCharTok{{-}}\FunctionTok{c}\NormalTok{(SAMPLE\_ID, Chromosome, Position), }\AttributeTok{names\_to =} \StringTok{"SNP\_ID"}\NormalTok{, }\AttributeTok{values\_to =} \StringTok{"Genotype"}\NormalTok{)}
\NormalTok{Teosinte\_long }\OtherTok{\textless{}{-}}\NormalTok{ Teosinte\_merge\_2 }\SpecialCharTok{\%\textgreater{}\%} \FunctionTok{mutate}\NormalTok{(}\FunctionTok{across}\NormalTok{(}\FunctionTok{everything}\NormalTok{(), as.character)) }\SpecialCharTok{\%\textgreater{}\%} 
  \FunctionTok{pivot\_longer}\NormalTok{(}\AttributeTok{cols =} \SpecialCharTok{{-}}\FunctionTok{c}\NormalTok{(SAMPLE\_ID, Chromosome, Position), }\AttributeTok{names\_to =} \StringTok{"SNP\_ID"}\NormalTok{, }\AttributeTok{values\_to =} \StringTok{"Genotype"}\NormalTok{)}
\end{Highlighting}
\end{Shaded}

\subsection{SNP Distribution by
Chromosome}\label{snp-distribution-by-chromosome}

\subsubsection{To compare the difference in the number of SNPs across
chromosomes for Maize and Teosinte I created a double bar
chart}\label{to-compare-the-difference-in-the-number-of-snps-across-chromosomes-for-maize-and-teosinte-i-created-a-double-bar-chart}

\subsubsection{I created a data table with the number of SNPs per
chromosome and then sorted that data table by chromosome
number}\label{i-created-a-data-table-with-the-number-of-snps-per-chromosome-and-then-sorted-that-data-table-by-chromosome-number}

\begin{Shaded}
\begin{Highlighting}[]
\NormalTok{Maize\_snp\_count }\OtherTok{\textless{}{-}}\NormalTok{ Maize\_long }\SpecialCharTok{\%\textgreater{}\%} \FunctionTok{count}\NormalTok{(Chromosome) }\SpecialCharTok{\%\textgreater{}\%} 
  \FunctionTok{mutate}\NormalTok{(}\AttributeTok{Chromosome =} \FunctionTok{factor}\NormalTok{(Chromosome, }\AttributeTok{levels =} \FunctionTok{c}\NormalTok{(}\DecValTok{1}\SpecialCharTok{:}\DecValTok{10}\NormalTok{, }\StringTok{"multiple"}\NormalTok{, }\StringTok{"unknown"}\NormalTok{))) }\SpecialCharTok{\%\textgreater{}\%} 
  \FunctionTok{arrange}\NormalTok{(Chromosome) }\SpecialCharTok{\%\textgreater{}\%} \FunctionTok{rename}\NormalTok{(}\StringTok{"SNPs"} \OtherTok{=}\NormalTok{ n)}
\NormalTok{Teosinte\_snp\_count }\OtherTok{\textless{}{-}}\NormalTok{ Teosinte\_long }\SpecialCharTok{\%\textgreater{}\%} \FunctionTok{count}\NormalTok{(Chromosome) }\SpecialCharTok{\%\textgreater{}\%} 
  \FunctionTok{mutate}\NormalTok{(}\AttributeTok{Chromosome =} \FunctionTok{factor}\NormalTok{(Chromosome, }\AttributeTok{levels =} \FunctionTok{c}\NormalTok{(}\DecValTok{1}\SpecialCharTok{:}\DecValTok{10}\NormalTok{, }\StringTok{"multiple"}\NormalTok{, }\StringTok{"unknown"}\NormalTok{))) }\SpecialCharTok{\%\textgreater{}\%} 
  \FunctionTok{arrange}\NormalTok{(Chromosome) }\SpecialCharTok{\%\textgreater{}\%} \FunctionTok{rename}\NormalTok{(}\StringTok{"SNPs"} \OtherTok{=}\NormalTok{ n)}
\end{Highlighting}
\end{Shaded}

\subsubsection{I combined the SNPs count data tables for Maize and
Teosinte by adding a species column. Then combined them using the rbind
command.}\label{i-combined-the-snps-count-data-tables-for-maize-and-teosinte-by-adding-a-species-column.-then-combined-them-using-the-rbind-command.}

\subsubsection{Then I again sorted them by chromosome
number}\label{then-i-again-sorted-them-by-chromosome-number}

\begin{Shaded}
\begin{Highlighting}[]
\NormalTok{Maize\_snp\_count}\SpecialCharTok{$}\NormalTok{Species }\OtherTok{\textless{}{-}} \StringTok{"Maize"}
\NormalTok{Teosinte\_snp\_count}\SpecialCharTok{$}\NormalTok{Species }\OtherTok{\textless{}{-}} \StringTok{"Teosinte"}
\NormalTok{Combined\_snp\_count }\OtherTok{\textless{}{-}} \FunctionTok{rbind}\NormalTok{(Maize\_snp\_count, Teosinte\_snp\_count)}
\NormalTok{Combined\_snp\_count}\SpecialCharTok{$}\NormalTok{Chromosome }\OtherTok{\textless{}{-}} \FunctionTok{factor}\NormalTok{(Combined\_snp\_count}\SpecialCharTok{$}\NormalTok{Chromosome, }\AttributeTok{levels =} \FunctionTok{c}\NormalTok{(}\DecValTok{1}\SpecialCharTok{:}\DecValTok{10}\NormalTok{, }\StringTok{"multiple"}\NormalTok{, }\StringTok{"unknown"}\NormalTok{))}
\end{Highlighting}
\end{Shaded}

\begin{Shaded}
\begin{Highlighting}[]
\FunctionTok{library}\NormalTok{(ggplot2) }\CommentTok{\# making sure ggplot2 is on}
\end{Highlighting}
\end{Shaded}

\subsubsection{I used ggplot to create the double bar chart. I used stat
= identity so that it uses the count that is already done. I also used
position = dodge to put the bar next to each other instead of
overlapping}\label{i-used-ggplot-to-create-the-double-bar-chart.-i-used-stat-identity-so-that-it-uses-the-count-that-is-already-done.-i-also-used-position-dodge-to-put-the-bar-next-to-each-other-instead-of-overlapping}

\subsubsection{\texorpdfstring{\textbf{With the bar chart you can see
that Maize has more SNPs per chromosome than Teosinte and the larger the
larger the chromosome number the less amount of SNPs
present}}{With the bar chart you can see that Maize has more SNPs per chromosome than Teosinte and the larger the larger the chromosome number the less amount of SNPs present}}\label{with-the-bar-chart-you-can-see-that-maize-has-more-snps-per-chromosome-than-teosinte-and-the-larger-the-larger-the-chromosome-number-the-less-amount-of-snps-present}

\begin{Shaded}
\begin{Highlighting}[]
\FunctionTok{ggplot}\NormalTok{(Combined\_snp\_count, }\FunctionTok{aes}\NormalTok{(}\AttributeTok{x =}\NormalTok{ Chromosome, }\AttributeTok{y =}\NormalTok{ SNPs, }\AttributeTok{fill =}\NormalTok{ Species)) }\SpecialCharTok{+} 
  \FunctionTok{geom\_bar}\NormalTok{(}\AttributeTok{stat =} \StringTok{"identity"}\NormalTok{, }\AttributeTok{position =} \StringTok{"dodge"}\NormalTok{) }\SpecialCharTok{+}
  \FunctionTok{labs}\NormalTok{(}\AttributeTok{title =} \StringTok{"SNP Distribution by Chromosome"}\NormalTok{, }\AttributeTok{x =} \StringTok{"Chromosome"}\NormalTok{, }\AttributeTok{y =} \StringTok{"Number of SNPs"}\NormalTok{)}
\end{Highlighting}
\end{Shaded}

\pandocbounded{\includegraphics[keepaspectratio]{R_assignment_Petersen_files/figure-latex/unnamed-chunk-27-1.pdf}}

\subsection{SNP Distribution per
Chromosome}\label{snp-distribution-per-chromosome}

\begin{Shaded}
\begin{Highlighting}[]
\NormalTok{Maize\_long}\SpecialCharTok{$}\NormalTok{Species }\OtherTok{\textless{}{-}} \StringTok{"Maize"}
\NormalTok{Teosinte\_long}\SpecialCharTok{$}\NormalTok{Species }\OtherTok{\textless{}{-}} \StringTok{"Teosinte"}
\NormalTok{Combined\_long }\OtherTok{\textless{}{-}} \FunctionTok{rbind}\NormalTok{(Maize\_long, Teosinte\_long)}
\NormalTok{Combined\_long}\SpecialCharTok{$}\NormalTok{Chromosome }\OtherTok{\textless{}{-}} \FunctionTok{factor}\NormalTok{(Combined\_long}\SpecialCharTok{$}\NormalTok{Chromosome, }\AttributeTok{levels =} \FunctionTok{c}\NormalTok{(}\DecValTok{1}\SpecialCharTok{:}\DecValTok{10}\NormalTok{, }\StringTok{"multiple"}\NormalTok{, }\StringTok{"unknown"}\NormalTok{))}
\NormalTok{Combined\_long }\OtherTok{\textless{}{-}}\NormalTok{ Combined\_long[Combined\_long}\SpecialCharTok{$}\NormalTok{Chromosome }\SpecialCharTok{\%in\%} \DecValTok{1}\SpecialCharTok{:}\DecValTok{10}\NormalTok{,] }\CommentTok{\# I took out the multiple and unknown chromosomes because they had no position data}
\end{Highlighting}
\end{Shaded}

\subsubsection{I needed to change the poisiton values to numerical data
in R so it could be
plotted}\label{i-needed-to-change-the-poisiton-values-to-numerical-data-in-r-so-it-could-be-plotted}

\begin{Shaded}
\begin{Highlighting}[]
\NormalTok{Combined\_long}\SpecialCharTok{$}\NormalTok{Position }\OtherTok{\textless{}{-}} \FunctionTok{as.numeric}\NormalTok{(}\FunctionTok{as.character}\NormalTok{(Combined\_long}\SpecialCharTok{$}\NormalTok{Position))}
\end{Highlighting}
\end{Shaded}

\begin{verbatim}
## Warning: NAs introduced by coercion
\end{verbatim}

\subsubsection{I used ploted both Maize and Teosinte on teh plots again
this time
overlapped}\label{i-used-ploted-both-maize-and-teosinte-on-teh-plots-again-this-time-overlapped}

\subsubsection{I used the facet wrap command to that all the chromosomes
had a different plot but where all still on the same axis for easy
comparison}\label{i-used-the-facet-wrap-command-to-that-all-the-chromosomes-had-a-different-plot-but-where-all-still-on-the-same-axis-for-easy-comparison}

\subsubsection{\texorpdfstring{\textbf{Again we see that maize has more
SNPs than Teosinte. Chromosome 7 is also interesting becuase there are
no SNPs in the middle of the
chromosome}}{Again we see that maize has more SNPs than Teosinte. Chromosome 7 is also interesting becuase there are no SNPs in the middle of the chromosome}}\label{again-we-see-that-maize-has-more-snps-than-teosinte.-chromosome-7-is-also-interesting-becuase-there-are-no-snps-in-the-middle-of-the-chromosome}

\begin{Shaded}
\begin{Highlighting}[]
\FunctionTok{ggplot}\NormalTok{(Combined\_long, }\FunctionTok{aes}\NormalTok{(}\AttributeTok{x =}\NormalTok{ Position, }\AttributeTok{fill =}\NormalTok{ Species)) }\SpecialCharTok{+}
  \FunctionTok{geom\_histogram}\NormalTok{(}\AttributeTok{binwidth =} \DecValTok{1000000}\NormalTok{, }\AttributeTok{position =} \StringTok{"identity"}\NormalTok{, }\AttributeTok{alpha =} \FloatTok{0.5}\NormalTok{) }\SpecialCharTok{+}
  \FunctionTok{facet\_wrap}\NormalTok{(}\SpecialCharTok{\textasciitilde{}}\NormalTok{Chromosome, }\AttributeTok{scales =} \StringTok{"free\_x"}\NormalTok{) }\SpecialCharTok{+}
  \FunctionTok{labs}\NormalTok{(}\AttributeTok{title =} \StringTok{"SNP Distribution per Chromosome"}\NormalTok{, }\AttributeTok{x =} \StringTok{"Position (bp)"}\NormalTok{, }\AttributeTok{y =} \StringTok{"Number of SNPs"}\NormalTok{)}
\end{Highlighting}
\end{Shaded}

\begin{verbatim}
## Warning: Removed 28050 rows containing non-finite outside the scale range
## (`stat_bin()`).
\end{verbatim}

\pandocbounded{\includegraphics[keepaspectratio]{R_assignment_Petersen_files/figure-latex/unnamed-chunk-30-1.pdf}}

\subsection{Missing data and Amount of
Heterozygosity}\label{missing-data-and-amount-of-heterozygosity}

\subsubsection{I created a new column that was labeled either Missing,
Homozygous, or Heterozygous using the ifelse logic
statements}\label{i-created-a-new-column-that-was-labeled-either-missing-homozygous-or-heterozygous-using-the-ifelse-logic-statements}

\begin{Shaded}
\begin{Highlighting}[]
\NormalTok{Combined\_long}\SpecialCharTok{$}\NormalTok{Type }\OtherTok{\textless{}{-}} \FunctionTok{ifelse}\NormalTok{(Combined\_long}\SpecialCharTok{$}\NormalTok{Genotype }\SpecialCharTok{\%in\%} \FunctionTok{c}\NormalTok{(}\StringTok{"?/?"}\NormalTok{), }\StringTok{"Missing"}\NormalTok{,}
                      \FunctionTok{ifelse}\NormalTok{(Combined\_long}\SpecialCharTok{$}\NormalTok{Genotype }\SpecialCharTok{\%in\%} \FunctionTok{c}\NormalTok{(}\StringTok{"A/A"}\NormalTok{, }\StringTok{"G/G"}\NormalTok{, }\StringTok{"T/T"}\NormalTok{, }\StringTok{"C/C"}\NormalTok{), }\StringTok{"Homozygous"}\NormalTok{,}
                      \StringTok{"Heterozygous"}\NormalTok{))}
\end{Highlighting}
\end{Shaded}

\subsubsection{I plotted the resulting data in a stacked bar chart to
compare Maize and
Teosinte.}\label{i-plotted-the-resulting-data-in-a-stacked-bar-chart-to-compare-maize-and-teosinte.}

\subsubsection{\texorpdfstring{\textbf{While the highest proportions of
SNPs came from homozygous genotypes, Teosinte has a higher proportion of
heterozygous genotypes than
Maize}}{While the highest proportions of SNPs came from homozygous genotypes, Teosinte has a higher proportion of heterozygous genotypes than Maize}}\label{while-the-highest-proportions-of-snps-came-from-homozygous-genotypes-teosinte-has-a-higher-proportion-of-heterozygous-genotypes-than-maize}

\begin{Shaded}
\begin{Highlighting}[]
\FunctionTok{ggplot}\NormalTok{(Combined\_long, }\FunctionTok{aes}\NormalTok{(}\AttributeTok{x =}\NormalTok{ Species, }\AttributeTok{fill =}\NormalTok{ Type)) }\SpecialCharTok{+}
  \FunctionTok{geom\_bar}\NormalTok{(}\AttributeTok{position =} \StringTok{"fill"}\NormalTok{) }\SpecialCharTok{+}
  \FunctionTok{labs}\NormalTok{(}\AttributeTok{y =} \StringTok{"Proportion"}\NormalTok{, }\AttributeTok{title =} \StringTok{"Proportion of Homozygous, Heterozygous, and Missing Sites"}\NormalTok{)}
\end{Highlighting}
\end{Shaded}

\pandocbounded{\includegraphics[keepaspectratio]{R_assignment_Petersen_files/figure-latex/unnamed-chunk-32-1.pdf}}

\subsection{Distribution of Genotype Type by
Chromosome}\label{distribution-of-genotype-type-by-chromosome}

\subsubsection{I plotted the proportion of each type of SNP
(heterozygous, homozygous, or Missing) for each chromosome for both
Maize and
Teosinte}\label{i-plotted-the-proportion-of-each-type-of-snp-heterozygous-homozygous-or-missing-for-each-chromosome-for-both-maize-and-teosinte}

\subsubsection{This will show is there are any chromosomes that are
extremely different in terms of the type of SNP for either Teosinte or
Maize}\label{this-will-show-is-there-are-any-chromosomes-that-are-extremely-different-in-terms-of-the-type-of-snp-for-either-teosinte-or-maize}

\subsubsection{\texorpdfstring{\textbf{Based on the plots there was not
any discernible trends that would warrant further
investigation}}{Based on the plots there was not any discernible trends that would warrant further investigation}}\label{based-on-the-plots-there-was-not-any-discernible-trends-that-would-warrant-further-investigation}

\begin{Shaded}
\begin{Highlighting}[]
\FunctionTok{ggplot}\NormalTok{(Combined\_long, }\FunctionTok{aes}\NormalTok{(}\AttributeTok{x =}\NormalTok{ Chromosome, }\AttributeTok{fill =}\NormalTok{ Type)) }\SpecialCharTok{+}
  \FunctionTok{geom\_bar}\NormalTok{(}\AttributeTok{position =} \StringTok{"fill"}\NormalTok{) }\SpecialCharTok{+}
  \FunctionTok{facet\_wrap}\NormalTok{(}\SpecialCharTok{\textasciitilde{}}\NormalTok{Species, }\AttributeTok{scales =} \StringTok{"free\_x"}\NormalTok{) }\SpecialCharTok{+} 
  \FunctionTok{labs}\NormalTok{(}\AttributeTok{title =} \StringTok{"Distribution of Genotype Types by Chromosome"}\NormalTok{,}
       \AttributeTok{x =} \StringTok{"Chromosome"}\NormalTok{,}
       \AttributeTok{y =} \StringTok{"Proportion"}\NormalTok{,}
       \AttributeTok{fill =} \StringTok{"Genotype Status"}\NormalTok{)}
\end{Highlighting}
\end{Shaded}

\pandocbounded{\includegraphics[keepaspectratio]{R_assignment_Petersen_files/figure-latex/unnamed-chunk-33-1.pdf}}

\section{END}\label{end}

\end{document}
